\chapter{Notes}

\textbf{Essence of Category Theory by Bartosz}
Category theory extracts the essence of structure and composition. At its foundation it deals with the composition of arrows. Building on composition of arrows it then goes on describing the ways objects can be composed: we have products, coproducts and, at a higher level, tensor products. They all describe various modes of composing objects. In monoidal categories any two objects can be composed.

So how do we approach problem solving? There is one general approach that we humans have developed for problem solving. We had to develop it because of the limitations of our brain, not because of the limitations of computers or our tools. Our brains have this relatively small cache memory, so when we’re dealing with a huge problem, we have to split it into smaller parts. We have to decompose bigger problems into smaller problems. And this is very human. This is what we do. We decompose, and then we attack each problem separately, find the solution; and once we have solutions to all the smaller problems, we recompose them.

So the essence of programming is composition.